\chapter{Лабораторная работа №4. Метод динамического программирования Беллмана}

\section*{Постановка (вариант 13)}
Объект: \(\dot x = A x + b u\), где
\[
A=\begin{bmatrix}0&1\\4&3\end{bmatrix},\quad b=\begin{bmatrix}2\\6\end{bmatrix},\quad Q=\begin{bmatrix}6&0\\0&3\end{bmatrix},\quad r=2,
\]
критерий на бесконечном интервале
\[
J=\int_0^{\infty} \big(x^T Q x + r u^2\big)\,dt,\qquad x(0)=\begin{bmatrix}1\\0\end{bmatrix}.
\]
Требуется синтезировать оптимальный регулятор методом Беллмана и промоделировать работу на заданном интервале времени.

\section{Уравнение Беллмана–Гамильтона–Якоби и CARE}
Полагаем, что оптимальная стоимость имеет квадратичный вид \(V(x)=x^T P x\) с \(P=P^T\succ0\). Тогда HJB
\[
0=\min_{u}\Big\{ x^T Q x + r u^2 + (\nabla V)^T (A x + b u) \Big\}
=\min_{u}\Big\{ x^T Q x + r u^2 + 2 x^T P (A x + b u) \Big\}.
\]
Минимум по \(u\) даёт \(u^*(x)= -\tfrac{1}{r} b^T P x\equiv -Kx\). Подстановка приводит к алгебраическому уравнению Риккати (CARE)
\[
A^T P + P A - P b r^{-1} b^T P + Q = 0,\qquad K=r^{-1} b^T P.
\]
При стабилизируемости \((A,b)\) и детектируемости \((Q^{1/2},A)\) существует единственное стабилизирующее \(P\succ0\), и \(J(\infty)=V(x_0)=x_0^T P x_0\).

% \subsection{Численная реализация}
% \begin{lstlisting}[caption={Bellman/LQR (вариант 13): CARE, интеграл стоимости и графики},label={lst:bellman13}]
% # lab4/python/bellman_var13.py
% import numpy as np
% from scipy.linalg import solve_continuous_are
% from scipy.integrate import solve_ivp
% import matplotlib.pyplot as plt

% A = np.array([[0.0, 1.0],[4.0, 3.0]])
% b = np.array([[2.0],[6.0]])
% Q = np.array([[6.0, 0.0],[0.0, 3.0]])
% r = 2.0

% P = solve_continuous_are(A, b, Q, r)
% K = (1.0/r) * (b.T @ P)
% Acl = A - b @ K

% x0 = np.array([1.0, 0.0])
% T = (0.0, 10.0)

% def cost(x, u):
%     return float(x.T @ Q @ x + r * u**2)

% def odefun(t, z):
%     x = z[:2]
%     u = -float(K @ x)
%     jdot = cost(x, u)
%     dx = Acl @ x
%     return np.hstack([dx, jdot])

% z0 = np.hstack([x0, 0.0])
% sol = solve_ivp(odefun, T, z0, max_step=0.01, rtol=1e-8, atol=1e-10)

% t = sol.t
% x = sol.y[:2,:]
% Jt = sol.y[2,:]
% u = -(K @ x).ravel()

% import os
% os.makedirs('/home/leonidas/projects/itmo/optimal-control-theory/lab4/images/task4', exist_ok=True)

% plt.figure(figsize=(6,4))
% plt.plot(t, x[0], label='x1'); plt.plot(t, x[1], label='x2')
% plt.xlabel('t'); plt.ylabel('states'); plt.grid(True, alpha=0.3)
% plt.legend(); plt.tight_layout()
% plt.savefig('/home/leonidas/projects/itmo/optimal-control-theory/lab4/images/task4/states.png', dpi=200)

% plt.figure(figsize=(6,4))
% plt.plot(t, u)
% plt.xlabel('t'); plt.ylabel('u'); plt.grid(True, alpha=0.3)
% plt.tight_layout()
% plt.savefig('/home/leonidas/projects/itmo/optimal-control-theory/lab4/images/task4/u.png', dpi=200)

% plt.figure(figsize=(6,4))
% plt.plot(t, Jt)
% plt.xlabel('t'); plt.ylabel('J(0,t)'); plt.grid(True, alpha=0.3)
% plt.tight_layout()
% plt.savefig('/home/leonidas/projects/itmo/optimal-control-theory/lab4/images/task4/J.png', dpi=200)

% print({'K': K.tolist(), 'J_final': float(Jt[-1]), 'x0TPx0': float(x0 @ (P @ x0))})
% \end{lstlisting}

\begin{figure}[H]
    \centering
    \begin{subfigure}{0.48\textwidth}
        \includegraphics{task4/states.png}
        \caption{Состояния $x_1,x_2$}
    \end{subfigure}\hfill
    \begin{subfigure}{0.48\textwidth}
        \includegraphics{task4/u.png}
        \caption{Управление $u(t)$}
    \end{subfigure}

    \vspace{0.5em}

    \begin{subfigure}{0.48\textwidth}
        \includegraphics{task4/J.png}
        \caption{Накопленный критерий $J(0,t)$}
    \end{subfigure}
    \caption{Метод Беллмана (вариант 13): моделирование}
\end{figure}

Получено: \(K\approx[1.1345\;1.8229]\). Инвариантная формула даёт \(J(\infty)=x_0^T P x_0\approx 2.4192\), что совпадает с интегральной оценкой по траектории.

\section{Отклонение коэффициентов регулятора}
Чтобы оценить устойчивость решения по критерию, умножим оба коэффициента оптимального \(K\) на \(0.8\) и \(1.2\). Итоговые значения приведены в таблице.
\begin{table}[H]
\centering
\begin{tabular}{lccc}
\toprule
Случай & \(k_1\) & \(k_2\) & \(J(\infty)=J(0,t\to\infty)\) \\
\midrule
Оптимальное \(K\) & 1.134 & 1.823 & 2.4192 \\
Уменьшено на \(20\%\) & 0.908 & 1.458 & 2.4364 \\
Увеличено на \(20\%\) & 1.361 & 2.187 & 2.4267 \\
\bottomrule
\end{tabular}
\caption{Изменение стоимости при небольшом отклонении \(K\)}
\end{table}
Во всех трёх случаях собственные значения замкнутой матрицы \(A-bK\) остаются в левой полуплоскости (\(-6.36\ldots-11.69\)), следовательно, устойчивость сохраняется. Отклонение на \(20\%\) даёт лишь \(0.3\dots0.7\%\) роста \(J(\infty)\), то есть критерий ухудшается совсем незначительно: регулятор демонстрирует хорошую робастность, что видно по кривым накопленного интеграла (рис.~\ref{fig:J_compare_pertK}).

\begin{figure}[H]
    \centering
    \includegraphics[width=0.65\textwidth]{task4/J_compare_pertK.png}
    \caption{Критерий \(J(0,t)\) при оптимальном и отклонённых значениях \(K\)}
    \label{fig:J_compare_pertK}
\end{figure}
